% Part: counterfactuals
% Chapter: minimal-change-semantics
% Section: transitivity

\documentclass[../../../include/open-logic-section]{subfiles}

\begin{document}

\olfileid{cnt}{min}{tra}

\olsection{સંક્રામકતા}

ભૌતિક શરતી વિધાન માટે શૃંખલા-નિયમ માન્ય છે: $!A \lif !B, !B
\lif !C \Entails !A \lif !C$. બીજા શબ્દોમાં, ભૌતિક શરતી વિધાન
સંક્રામક છે. શું પ્રતિતથ્યાત્મક વિધાનો માટે પણ એવું જ છે? સ્ટાલનેકરે
આપેલું નીચેનું ઉદાહરણ જુઓ.
\begin{quote}
  જો જે.~એડગર હૂવર રશિયન તરીકે જન્મ્યા હોત, તો તેઓ કોમ્યુનિસ્ટ બન્યા હોત.

  જો જે.~એડગર હૂવર કોમ્યુનિસ્ટ હોત, તો તેઓ દેશદ્રોહી હોત.

  તેથી, જો જે.~એડગર હૂવર રશિયન તરીકે જન્મ્યા હોત, તો તેઓ દેશદ્રોહી હોત.
\end{quote}
હૂવર વાસ્તવમાં જે સમયે જન્મ્યા એ જ સમયે અમેરિકાને બદલે રશિયામાં
જન્મ્યા હોત, તો તેઓ સોવિયેત સંઘમાં ઊછરીને કોમ્યુનિસ્ટ બન્યા હોત
(એવું ધારી લઈએ). તેથી પહેલો પૂર્વાધાર સત્ય છે. એ જ રીતે, બીજા
પૂર્વાધારને અલગથી લઈએ તો તે પણ સત્ય છે. પરંતુ તારણ અસત્ય છે:
હૂવર સોવિયેત સંઘમાં જન્મ્યા હોત, તો સંભવતઃ તેઓ કટ્ટર કોમ્યુનિસ્ટ
બન્યા હોત અને (પોતાના દેશના) દેશદ્રોહી ન બન્યા હોત. સત્યમૂલ્યોનું
આ સહજ નિયોજન સ્ટાલનેકર--લૂઈસના વિશ્લેષણથી પણ સમર્થિત થાય છે.
હૂવરનું જન્મસ્થળ બદલાય એટલો જ ફેરફાર ધરાવતું આપણા જગતથી સૌથી
નજીકનું શક્ય જગત એ છે જેમાં હૂવર સોવિયેત સંઘના સારા નાગરિક તરીકે
ઊછરે છે. પહેલો પૂર્વાધાર અને તારણ, બંનેનો પૂર્વપક્ષ એ જગતમાં સત્ય
છે; અને ત્યાં હૂવર કોમ્યુનિસ્ટ પક્ષના વફાદાર સભ્ય છે, તેથી દેશદ્રોહી
નથી. પરંતુ બીજો પૂર્વાધાર મૂલવવા જુદું જગત જોવું પડે: હૂવર
કોમ્યુનિસ્ટ હોય એવું સૌથી નજીકનું જગત. એમાં તેઓ અમેરિકામાં જન્મ્યા,
પછી કોમ્યુનિસ્ટ બન્યા અને તેથી દેશદ્રોહી બન્યા.\footnote{અલબત્ત, આ
  ઉદાહરણનું બળ સમજવા કેટલીક અધિભૌતિક અને રાજકીય ધારણાઓ સ્વીકારવી
  પડે; દાખલા તરીકે, હૂવર રશિયન માતાપિતાને જન્મ્યા હોત એ શક્યતા,
  અથવા 1950ના દાયકામાં અમેરિકાના કોમ્યુનિસ્ટો પોતાના દેશના
  દેશદ્રોહી હતા એ ધારણા.}

\begin{prob}
  પ્રતિતથ્યાત્મક વિધાનોની સંક્રામકતા નિષ્ફળ જાય તેનું વિશ્વસનીય,
  સહજ ઉદાહરણ શોધો.
\end{prob}

\begin{ex}\ollabel{ex:trans-counterex}
  ગોળા-અર્થવિચાર આ અનુમાનને અમાન્ય ઠરાવે છે; એટલે $p
  \cif q, q \cif r \Entails/ p \cif r$. નિદર્શ $\mModel{M}
  = \tuple{W, O, V}$ લો જેમાં $W = \{w, w_1, w_2\}$,
  $O_w = \{\{w\}, \{w, w_1\}, \{w, w_1, w_2\}\}$,
  $V(p) = \{w_2\}$, $V(q) = \{w_1, w_2\}$ અને
  $V(r) = \{w_1\}$. $p$ને સ્વીકારતો ગોળો $S = \{w, w_1,
  w_2\}$ છે અને તેમાંનાં બધાં જગતોમાં $p \lif q$ સત્ય છે; તેથી
  $\mSat{M}{p \cif q}[w]$. વળી $q$ને સ્વીકારતો ગોળો
  $S' = \{w, w_1\}$ છે અને તેમાંનાં બધાં જગતોમાં
  $q \lif r$ સત્ય છે; તેથી
  $\mSat{M}{q \cif r}[w]$. પરંતુ $p$ને સ્વીકારતા ગોળા
  $\{w, w_1, w_2\}$માં જગત~$w_2$ છે, જ્યાં
  $\mSat/{M}{p \lif r}[w_2]$.
  \sourcecorrection{OLMD-110}{મૂળમાં q→r બધા જગતોમાં સત્ય કહેતાં ખોટું અસંતોષ-સૂત્ર મૂકાયું છે; અહીં તે જ ભૌતિક શરતી વિધાન બધા જગતોમાં સત્ય હોવાનું સ્પષ્ટ લખ્યું છે.}
\end{ex}

\begin{prob}
\olref[cnt][min][tra]{ex:trans-counterex}ના પ્રત્યુદાહરણને અનુરૂપ
ગોળાઓનો આલેખ દોરો.
\end{prob}

\begin{prob}
  \olref[cnt][min][tra]{ex:trans-counterex}માં જગત $w_2$ પર હૂવર
  રશિયામાં જન્મ્યા છે, કોમ્યુનિસ્ટ છે અને દેશદ્રોહી નથી; જગત $w_1$
  પર તેઓ અમેરિકામાં જન્મ્યા છે, કોમ્યુનિસ્ટ છે અને દેશદ્રોહી છે.
  આ નિદર્શમાં $w_1$,~$w$ની $w_2$ કરતાં વધુ નજીક છે. શું આ જરૂરી
  છે? રશિયામાં જન્મવું, હૂવર કોમ્યુનિસ્ટ બને તેની તુલનામાં વધુ દૂરની
  શક્યતા છે એવી ધારણા વગરનું પ્રત્યુદાહરણ આપી શકો છો?
\end{prob}

\end{document}
